\documentclass[10pt,a4paper,twocolumn,twoside]{article}
\usepackage[utf8]{inputenc}
\usepackage[catalan]{babel}
\usepackage{multicol}
\usepackage{graphicx}
\usepackage{fancyhdr}
\usepackage{times}
\usepackage{titlesec}
\usepackage{multirow}
\usepackage{lettrine}
\usepackage[top=2cm, bottom=1.5cm, left=2cm, right=2cm]{geometry}
\usepackage[figurename=Fig.,tablename=TAULA]{caption}
\captionsetup[table]{textfont=sc}

\author{\LARGE\sffamily Pablo Chen}
\title{\Huge{\sffamily Disseny i implementació d'un joc de puzle tipus Sudoku per a aprendre els números en braille incorporant dispositius d'interacció gestual}}
\date{}

\newcommand\blfootnote[1]{%
  \begingroup
  \renewcommand\thefootnote{}\footnote{#1}%
  \addtocounter{footnote}{-1}%
  \endgroup
}

\titleformat{\section}
{\large\sffamily\scshape\bfseries}
{\textbf{\thesection}}{1em}{}

\begin{document}

\fancyhead[LO]{\scriptsize Pablo Chen: Disseny i implementació d'un joc de puzle tipus Sudoku per a aprendre els números en braille}
\fancyhead[RO]{\thepage}
\fancyhead[LE]{\thepage}
\fancyhead[RE]{\scriptsize EE/UAB TFG INFORMÀTICA:\@Disseny i implementació d'un Sudoku en braille}

\fancyfoot[CO,CE]{}

\fancypagestyle{primerapagina}
{
  \fancyhf{}
  \fancyhead[L]{\scriptsize TFG EN ENGINYERIA INFORMÀTICA, ESCOLA D'ENGINYERIA (EE), UNIVERSITAT AUTÒNOMA DE BARCELONA (UAB)}
  \fancyfoot[C]{\scriptsize juny de 2022, Escola d'Enginyeria (UAB)}
}

\renewcommand{\headrulewidth}{0pt}
\renewcommand{\footrulewidth}{0pt}
\pagestyle{fancy}

\twocolumn[\begin{@twocolumnfalse}

    \maketitle

    \thispagestyle{primerapagina}
    \begin{center}
      \parbox{0.915\textwidth} {\sffamily \textbf{Resum--} Resum del projecte, màxim 10
        línies.\\Resum del projecte, màxim 10 línies.\\Resum del projecte, màxim 10
        línies.\\Resum del projecte, màxim 10 línies.\\Resum del projecte, màxim 10
        línies.\\Resum del projecte, màxim 10 línies.\\Resum del projecte, màxim 10
        línies.\\Resum del projecte, màxim 10 línies.\\Resum del projecte, màxim 10
        línies.\\Resum del projecte, màxim 10 línies.\\Resum del projecte, màxim 10
        línies.\\Resum del projecte, màxim 10 línies.\\Resum del projecte, màxim 10
        línies.\\Resum del projecte, màxim 10 línies.\\ \textbf{Paraules clau-- }
        Paraules clau del treball, màxim 2 línies\\ \\ \bigskip \\ \textbf{Abstract--}
        Versió en anglès del resum Resum del projecte, màxim 10 línies.\\Resum del
        projecte, màxim 10 línies.\\Resum del projecte, màxim 10 línies.\\Resum del
        projecte, màxim 10 línies.\\Resum del projecte, màxim 10 línies.\\Resum del
        projecte, màxim 10 línies.\\Resum del projecte, màxim 10 línies.\\Resum del
        projecte, màxim 10 línies.\\Resum del projecte, màxim 10 línies.\\Resum del
        projecte, màxim 10 línies.\\Resum del projecte, màxim 10 línies.\\ projecte,
        màxim 10 línies.\\Resum del projecte, màxim 10 línies.\\ \\ \\
        \textbf{Keywords-- } Versió en anglès de les paraules clau.
        \\}

      \bigskip

      {\vrule depth 0pt height 0.5pt width 4cm\hspace{7.5pt}%
        \raisebox{-3.5pt}{\fontfamily{pzd}\fontencoding{U}\fontseries{m}\fontshape{n}\fontsize{11}{12}\selectfont\char70}%
        \hspace{7.5pt}\vrule depth 0pt height 0.5pt width 4cm\relax}

    \end{center}

    \bigskip
    %\end{abstract}
  \end{@twocolumnfalse}]

\blfootnote{$\bullet$ E-mail de contacte: pablo.chen@uab.cat}
\blfootnote{$\bullet$ Menció realitzada: Enginyeria de Computadors / Computació / Enginyeria del Software / Tecnologies de la Informació}
\blfootnote{$\bullet$ Treball tutoritzat per: Enric Martí Godia (Computació)}
\blfootnote{$\bullet$ Curs 2021/22}

\section{Introducció}\label{sec:introduction}

\lettrine[lines=3] El Sudoku és un trencaclosques on el jugador tracta d'omplir una graella
parcialment emplenada de 9 $\times$ 9 caselles dividida en subgraelles de 3
$\times$ 3 (anomenades regions o blocs) amb un set de números de manera que
cada fila, columna i subgraella contingui els números només una vegada. El
sudoku modern va ser popularitzat en el 2004 quan Wayne Gould va desenvolupar
un programa informàtic per produir trencaclosques únics ràpidament.

Per una altra banda, el braille és un sistema d'escriptura i lectura tàctil
basat en 6 punts disposats en una graella de 3 $\times$ 3 utilitzat per a
persones cegues o amb mala visió.

% TODO EXPLICAR BRAILE

\section{Estat de l'art}\label{sec:stateOftheArt}
% TODO EXPLICAR BRAILE

\section{Objectiu}\label{sec:objective}
L'objectiu principal del projecte consisteix en desenvolupar un sudoku en 3D
des de zero. El joc tracta d'una graella de 9 $\times$ 9 fitxes on es
visualitzarà el valor numèric amb l'alfabet braille. Per una altra banda,
mitjançant el teclat i el ratolí, el jugador podrà interactuar amb les caselles
i les fitxes per veure el valor numèric de la fitxa. Concretament, es volen
aconseguir els següents objectius:
\begin{itemize}
  \item Desenvolupar un model en 3D del sudoku.
  \item Incorporar ajudes per aprendre els números o l'alfabet en braille.
\end{itemize}

Per aconseguir-lo, es desenvoluparà un petit conjunt de classes que permetrà
renderitzar objectes en la pantalla i tractar els esdeveniments dels
perifèrics.

El joc haurà de gestionar els següents esdeveniments dels perifèrics i
renderitzar les següents formes geomètriques:
\begin{itemize}
  \item Pulsacions del teclat.
  \item Pulsacions del ratolí, un clic.
  \item Moviments del ratolí.
  \item Moviments de la roda del ratolí.
  \item Moviments de la finestra del joc.
  \item Canvi de la finestra del joc.
  \item Canvis d'estats de la finestra.
  \item Un cub.
  \item Una esfera.
  \item Un pla.
\end{itemize}

El joc haurà d'implementar les següents funcionalitats:
\begin{itemize}
  \item Guardar i carregar partides a partir d'un fitxer de text.
  \item Generar un sudoku aleatori d'una determinada dificultat.
  \item Donar pistes en la resolució del sudoku de manera interactiva.
  \item Aconseguir l'execució en diverses plataformes, principalment Windows i Linux.
\end{itemize}
% Resumen: presente (también pasado)
% Introducción: presente
% Materiales, metodología, procedimientos: pasado
% Resultados: pasado (o presente)
% Conclusiones: presente

\section{Metodologia}\label{sec:methodology}
A continuació s'expliquen la metodologia seguida i les eines utilitzades al
llarg d'aquest projecte.

\subsection{Kanban}
Per a poder optimitzar al màxim els recursos necessaris, s'ha seguit una
metodologia àgil, celebrant revisions periòdiques generals en finalitzar cada
sprint on es van veure els resultats de la fase anterior, es van valorar els
avenços i es va planificar el següent sprint, concretament, s'ha utilitzat una
tècnica anomenada Kanban, es van dividir les tasques en diferents subtasques i
es van col·locar cada tasca segons el progrés (to do, in progress, done). En
particular, s'ha fet servir KanbanFlow per gestionar les subtasques.

Finalment, per a gestionar el codi font, s'ha creat un repositori a GitHub on
s'ha anat pujant a ells els avenços fets.

\subsection{Eines}

Respecte a les eines i llibreries de desenvolupament, s'ha utilitzat
principalment, l'entorn de desenvolupament multiplataforma CLion per C/C++ per
poder desenvolupar el joc tant en Windows com en Linux, Renderdoc per a
diagnosticar problemes de renderització. Respecte a les llibreries, s'ha fet
servir ImGUI per generar la interfície d'usuari, GLFW per gestionar la
inicialització de la finestra i obtenir els estats dels perifèrics, OpenGL per
interactuar amb la GPU, la llibreria glbinding per carregar les funcions
d'OpenGL i, finalment, GLM per realitzar els càlculs amb matrius i vectors.

\section{Desenvolupament}\label{sec:development}

\subsection{Classes de renderització}
\subsubsection{Cub}
\subsubsection{Esfera}
\subsubsection{Pla}
\subsubsection{Punter}
\subsection{Classes d'esdeveniments}
\subsection{Classe càmera}
\subsection{Entitat principal}
\subsection{Sudoku}
\subsection{Classe de control de l'estat}
\subsection{Procediments auxiliars}

\section{Resultats}\label{sec:results}
\section{Conclusions}\label{sec:conclusion}

\section*{Agraïments}

\begin{thebibliography}{11}
  \bibitem{latex}
  http://en.wikibooks.org/wiki/LaTeX

  \bibitem{2}
  Referència 2

  \bibitem{3}
  Etc.

\end{thebibliography}

\appendix

\section*{Apèndix}

\setcounter{section}{1}

\subsection{Secció d'Apèndix}

\end{document}